\begin{abstract}

	山区洪涝灾害常导致道路中断、通信基站损毁，无人机成为打通“最后一公里”应急物资投送与通信保障的关键装备。本文围绕\textbf{临时调度中心 O01 与 15 个服务区、80 个不可拆货箱、8 架异构运输无人机、2 架中继无人机}构成的协同任务，建立“\textbf{统一物理口径—四问递进建模—独立校验器复核}”的模型体系，并对若干违反直觉的结论给出定量论证。全文严格区分三类信息：\textbf{附件直接核对的输入}、\textbf{本文程序复现并经验证的结果}、以及\textbf{仅作对照的候选结果}。

	\textbf{对于问题一}，先由 30\,m DEM 沿航段采样最高地面高程，按“沿线最高点 $+50$\,m”确定巡航海拔，逐段求得 240 条航段的水平距离、爬升与下降高度（实测单向距离 \SIrange{2.808}{8.079}{km}、巡航海拔 \SIrange{271}{586}{m}）。据此建立\textbf{载荷—航程关系} $L_{g}(q)=L_{g0}-(L_{g0}-L_{gF})(q/\Qg)^{3/2}$ 与逐段能耗模型：水平巡航能耗由航程反推 $E^{hor}=(d/\Lg(q))\Euse$，爬升附加能耗按效率折算 $E^{up}=mgh^{+}/\eta_{up}$，下降能耗效率为 0 故不单独计。最大安全载荷由能量约束\textbf{取等号}定义；因指数 $3/2$ 无初等反函数，用\textbf{单调性 $+$ Brent 法数值反解}并强制回代，实测相对误差 $5.6\times10^{-15}$，达机器精度。结果表明\textbf{结构载重而非能量才是主导约束}：A 型 $15/15$ 个服务区全部受 $\Qg$ 约束，B 型 $14/15$，仅 C 型在 5 个远距离服务区受能量约束（S008 反解得 \SI{58.99}{kg}）。随后把每个服务区化为\textbf{质量—体积二维装箱}问题用 FFD 求解，并推导 \textbf{Martello--Toth 型解析下界}。最终得到 \textbf{18 个架次}（全部 C 型）的组批方案，\textbf{逐服务区等于解析下界，架次数维度可证最优}；总运输能耗 \SI{75.07}{kWh}，累计作业 \SI{9.37}{h}。$\rhog$ 敏感性显示：$\rhog$ 由 $0.20$ 增至 $0.35$ 使架次由 18 增至 25、能耗升至 \SI{106.60}{kWh}；$\rhog\ge0.375$ 时部分远距离服务区\textbf{即使空载也无法返航，问题开始无可行解}。

	\textbf{对于问题二}，将单点往返扩展为\textbf{异构多点多架次调度}，联合决策货箱组批、访问顺序、机型、实体无人机、共享电池与各架次开始时刻。多点串飞的载荷沿航段递减、能耗必须\textbf{逐段累加}（不能把多段距离相加当作一段），交接时间按 $\sum_{\text{站}}(t_{h0}+k_st_{h1})$ \textbf{逐站累加}。采用“构造—选序—局部搜索—时段驱动贪心派发”四阶段求解，并以\textbf{两阶段充电模型}刻画电池周转；其中首批箱按\textbf{机队槽位轮转播种}（4 架 A $\to$ 2 架 B $\to$ 2 架 C），派发平局按\textbf{最小松弛量} $\text{slack}=\min_b(d_b-t_b^{deliver})$ 而非能耗决出。得到 \textbf{25 个架次}（A 型 8 $+$ B 型 10 $+$ C 型 7，8 架无人机全部投入使用）、总能耗 \SI{77.30}{kWh}、完工 \SI{3.02}{h}（\SI{10855.1}{s}）、期望送达准时率 \textbf{100.0\%}（80/80 箱全部按时送达），首批违规与期望违规均为 0 箱。需要特别说明：\textbf{压缩架次数的正确做法不是调高架次权重}——“权重法”虽能把架次数压到 25，却让真实排程冒出 4 箱超期（准时率降至 95.0\%）而目标函数看不见；本文改为以“\textbf{真实调度器判定 0 违规}”为硬前提、以“架次数最少”为唯一方向，对每一次两两合并都过一遍真实调度器、不通过即回退。25 个架次距理论下界 24（$3$ 轮 $\times$ 8 架）仅高 1 个，质量装载率 \textbf{79\%}。

	\textbf{对于问题三}，新增\textbf{连续通信约束}：运输无人机在爬升、巡航、下降与投送全程不得中断。沿完整轨迹以 $\Delta t=2$\,s 逐时刻采样并作“直连 $>$ 中继 $>$ 中断”三态判定，实测 25 个运输架次中 \textbf{20 个存在直连中断}（平均中断占比 28.4\%），证明\textbf{中继为必需项}。链路预算给出三条双向门限：直连 \SI{122}{dB}、中继接入 \SI{116}{dB}、中继回传 \SI{126}{dB}，\textbf{接入段最弱}，故中继应\textbf{贴近作业空域而非贴近网关}。本文把“连续轨迹 $\times$ 连续三维选址”降维为可验证的充分条件，并特别强调\textbf{覆盖率有两个不可混用的口径}：\textbf{几何可达覆盖率}只回答“是否存在合适的悬停点”，而\textbf{时间轴覆盖率}才回答“中继那一刻是否在站”。题目仅配置 2 架中继，其可用能量 \SI{2.56}{kWh} 除以服务功率 \SI{1.10}{kW} 决定\textbf{单次在站时长上限约 \SI{140}{min}}，而 20 个架次的\textbf{真实中断时长合计约 \SI{168.5}{min}、最大并发 6 架次}，故按时间轴复核\textbf{无法保障全部架次}。本文按“服务窗口为硬约束”重排（中继须在每个中断区间起点前完成建链），给出 \textbf{6 个中继架次}、\textbf{时间轴覆盖率 $11/20=55.0\%$}，并如实报告\textbf{中继资源缺口 4 架}；运输与中继总能耗 \SI{79.65}{kWh}，联合完工 \SI{3.05}{h}。

	\textbf{对于问题四}，把“同一架次涉及的服务区必须划入同一任务组”形式化为服务区图上的\textbf{等价关系}，任务分区即求\textbf{连通分量（原子单元）}。对问题三的 25 个架次构图，其中 12 个多点架次把 15 个服务区串成 \textbf{3 个连通分量}，故\textbf{合法组数 $K$ 只能取 1 至 3}；\textbf{$K=2$ 与 $K=3$ 均可由原子单元直接合并得到，改动 0 个架次}。资源核算表明 $K=1/2/3$ 的四类资源总量为 \textbf{33/34/38}（台$\cdot$组）、组间不均衡度 $0/1.8517/2.7035$、单组最大工作量由 \SI{13.95}{h} 降至 \SI{13.00}{h}，即\textbf{分区不节省资源却能降低单组峰值工作量}；最关键的缺口在中继（库存 2 架 vs 需 4$\sim$5 架），三种分区下都存在且\textbf{分区无法缓解}。

	本文的主要特色在于：\textbf{①物理口径唯一}——附录给定的时间、能耗、充电、链路公式只在公共层实现一次，四问共用，杜绝口径分叉；\textbf{②结果可独立复算}——独立校验器\textbf{不导入任何求解器模块}，从附件与 DEM 重算全部物理量，用 18 类约束（含负样本测试）复核，问题二\textbf{违规 0 条}、问题三\textbf{硬违规 0 条}；\textbf{③敢于纠正否定性结论}——早期版本曾因“机队坍缩 $+$ 派发退化”误判“首批保障时限物理不可行”，又曾因“中继排班不认服务窗口 $+$ 校验器从未执行该项检查”误报“通信覆盖 100\%”；本文均给出根因与修正，并\textbf{不把几何可达当作时间轴已保障上报}。

	\keywords{异构多架次调度；集合划分；二维装箱；解析下界；连续通信约束；中继选址；图连通分量；独立可行性校验}
\end{abstract}
